\appendix

\chapter{Prompts, Judge Instructions and Atom Checklists}
\label{app:prompts}

This appendix reproduces the generation prompts, the atomic MQM
judge prompt, and a representative atom checklist used in the
experiments of Chapter~\ref{chap:implementation}.

\section{Generation Prompts (English)}
\label{app:prompts-en}

\subsection{Line Plot}

\begin{codebox}[title={English line-plot generation prompt}]
You are an annotator tasked with describing scientific line
plots in a structured but natural paragraph format.

Your goal is to clearly describe what is shown, without making
inferences or drawing conclusions. Stick to objective features
only. Use neutral, technical language.

For each plot, write a single paragraph that includes:

  A description of the overall purpose of the plot (what is
  being illustrated, compared, or represented)

  A description of the x-axis, including its label, scale type
  (e.g., linear or logarithmic), value range, units (if
  present), and tick label interval if numeric and regularly
  spaced

  A description of the y-axis, following the same structure as
  the x-axis

  A description of each line in the plot, including its color,
  marker shape, line style (e.g., solid, dashed), and how it
  changes across the x-axis. Where possible, include
  representative values at the start and end of the x-axis,
  and note any peaks, declines, or steady trends using exact
  data points from the plot (not inferred ones)

Avoid interpreting the cause of changes or making comparisons
between lines. Do not include section titles; write everything
in flowing sentences within a single paragraph per plot.
\end{codebox}

\subsection{Bar Chart}

\begin{codebox}[title={English bar-chart generation prompt}]
You are an annotator tasked with describing scientific bar
charts in a structured but natural paragraph format.

Your goal is to objectively describe what is shown, using
neutral and technical language. Avoid making inferences,
comparisons, or drawing conclusions about the data.

For each bar chart, write a single paragraph that includes:

  A description of the overall purpose of the chart (what is
  being illustrated, compared, or represented)

  A description of the axes: indicate which axis contains the
  categories (horizontal or vertical), and describe the axis
  label, order, and scale type if present (e.g., linear,
  logarithmic). Then describe the value axis, including its
  label, scale, value range, units (if present), and tick
  label interval if numeric and regularly spaced

  A description of each bar (or grouped/stacked set of bars),
  including the label, color, and approximate value or length.
  Where possible, include representative values from the
  shortest and tallest bars or note consistent trends in
  height or length, using exact data from the chart

  If the chart contains grouped or stacked bars, explain how
  the grouping or stacking is visually arranged and what each
  component represents based on labels or legends

  Mention whether the bars are sorted in a specific way (e.g.,
  descending, alphabetical), and whether any bars are visually
  emphasized (e.g., shaded, bolded, outlined, or annotated)

Avoid interpreting causes or trends, and do not compare
categories unless the comparison is explicitly labeled. Focus
solely on visual and labeled characteristics.

Write everything in a single, continuous paragraph without
section headers or bullet points.
\end{codebox}

\subsection{Pie Chart}

\begin{codebox}[title={English pie-chart generation prompt}]
You are an annotator tasked with describing scientific pie
charts in a structured but natural paragraph format.

Your goal is to objectively describe what is shown in the
chart, using neutral and technical language. Avoid interpreting
or explaining the significance of the proportions.

For each chart, write a single paragraph that includes:

  A description of the overall purpose of the chart (what is
  being represented or compared)

  The total number of slices (or categories) shown

  For each slice, include the label, color (if distinguishable),
  and approximate percentage or proportion if shown or easily
  estimated

  Mention whether the chart includes a legend, percentage
  labels inside the slices, or if any slices are visually
  emphasized (e.g. exploded, offset, or bolded)

  If the segments appear ordered in a particular way (e.g.
  descending size), note this as a visual observation, not
  interpretation

Write everything in a single, natural paragraph. Do not include
section titles or bullet points. Avoid drawing conclusions or
comparing slices beyond what is visually and explicitly
presented.
\end{codebox}


\section{Generation Prompt (German Translation)}
\label{app:prompts-de}

\begin{codebox}[title={German bar-chart generation prompt (translated)}]
Sie sind ein Annotator, dessen Aufgabe es ist,
wissenschaftliche Balkendiagramme in einem strukturierten,
aber natuerlichen Absatzformat zu beschreiben.

Ihr Ziel ist es, das Dargestellte objektiv und in einer
neutralen, fachlichen Sprache zu beschreiben. Vermeiden Sie
es, Schlussfolgerungen zu ziehen, Vergleiche anzustellen
oder Rueckschluesse auf die Daten zu ziehen.

Verfassen Sie fuer jedes Balkendiagramm einen einzelnen
Absatz, der Folgendes enthaelt:

  Eine Beschreibung des Gesamtzwecks des Diagramms (was
  dargestellt, verglichen oder repraesentiert wird)

  Eine Beschreibung der Achsen: Geben Sie an, welche Achse
  die Kategorien enthaelt (horizontal oder vertikal), und
  beschreiben Sie die Achsenbeschriftung, die Reihenfolge
  und den Skalentyp, falls vorhanden (z. B. linear,
  logarithmisch). Beschreiben Sie dann die Werteachse,
  einschliesslich ihrer Beschriftung, Skala, Wertebereich,
  Einheiten (falls vorhanden) und Tick-Beschriftungsintervall

  Eine Beschreibung jedes Balkens (oder jeder gruppierten/
  gestapelten Gruppe von Balken), einschliesslich
  Beschriftung, Farbe und ungefaehrer Wert oder Laenge

  Wenn das Diagramm gruppierte oder gestapelte Balken
  enthaelt, erklaeren Sie, wie die Gruppierung oder
  Stapelung visuell angeordnet ist

  Geben Sie an, ob die Balken in einer bestimmten
  Reihenfolge sortiert sind und ob bestimmte Balken visuell
  hervorgehoben sind

Schreiben Sie alles in einem einzigen, zusammenhaengenden
Absatz ohne Abschnittsueberschriften oder Aufzaehlungspunkte.
\end{codebox}


\section{Atomic MQM Judge Prompt}
\label{app:judge-prompt}

\begin{codebox}[title={Atomic MQM judge system prompt (excerpt)}]
You are an expert evaluator assessing the quality of a
machine-generated scientific figure description using an
atomic checklist approach.

You are given the FIGURE IMAGE, an ATOM CHECKLIST of
verifiable facts extracted from a human expert annotation,
the full REFERENCE DESCRIPTION, and a MACHINE-GENERATED
description.

Step 1 - Accuracy Check (atoms + image):
  For each atom, check if the model description mentions it
  AND gets it right. If wrong, flag an Accuracy error.

Step 2 - Completeness Check (atoms):
  For each atom, check if the model description covers it
  AT ALL. If missing, flag a Completeness error.

Step 3 - Hallucination Check (atoms + image):
  Flag claims NOT covered by any atom that are factually
  incorrect.

Step 4 - Clarity Check (reference description):
  Compare against the reference for readability.

Severity assignment: atom severity caps error severity.
  Critical/Important atom, completely wrong -> Major
  Critical/Important atom, partially off    -> Minor
  Minor atom, any degree of error           -> Minor

Output: {"errors": [{"category": "...", "sub_type": "...",
  "severity": "...", "text_span": "...",
  "evidence": "...", "atom_id": "..."}]}
\end{codebox}


\section{Atom Checklist: \texttt{german\_fig\_023}}
\label{app:atoms}

\begin{codebox}[title={Complete atom checklist for \texttt{german\_fig\_023} (26 atoms)}]
{ "figure_key": "german_fig_023",
  "figure_type": "bar_chart",
  "language": "German",
  "atoms": [
    { "id": "chart_type",     "severity": "critical",
      "value": "Die Abbildung besteht aus zwei Diagrammen" },
    { "id": "chart_purpose",  "severity": "critical",
      "value": "zeigt die Investitionen der
       Gebietskoerperschaften in Relation zum BIP" },
    { "id": "left_diagram",   "severity": "critical",
      "value": "Das linke Diagramm stellt die
       Bruttoanlageinvestitionen" },
    { "id": "right_diagram",  "severity": "critical",
      "value": "das rechte Diagramm zeigt die
       Nettoanlageinvestitionen" },
    { "id": "x_axis_range",   "severity": "critical",
      "value": "Die horizontale Achse umfasst die Jahre
       1991 bis 2023" },
    { "id": "y_axis_unit",    "severity": "critical",
      "value": "die vertikale Achse ist in Prozent des
       BIP skaliert" },
    { "id": "y_axis_scale",   "severity": "important",
      "value": "reicht von -0,5 bis 3,5" },
    { "id": "bar_bund",       "severity": "critical",
      "value": "\"Bund\" (blau)" },
    { "id": "bar_laender",    "severity": "critical",
      "value": "\"Laender\" (rot)" },
    { "id": "bar_gemeinden",  "severity": "critical",
      "value": "\"Gemeinden\" (gelb)" },
    { "id": "bar_gesamt",     "severity": "critical",
      "value": "\"Gesamt\" als schwarze Linie" },
    { "id": "brutto_trend",   "severity": "important",
      "value": "Die Bruttoanlageinvestitionen zeigen
       einen leichten Rueckgang" },
    { "id": "netto_trend",    "severity": "important",
      "value": "Die Nettoanlageinvestitionen bewegen sich
       teilweise im negativen Bereich" },
    { "id": "bund_brutto",    "severity": "important",
      "value": "Der Bund traegt einen relativ kleinen
       Anteil bei" },
    { "id": "gemeinden_brutto","severity": "important",
      "value": "Die Gemeinden haben den groessten Anteil
       an den Bruttoinvestitionen" },
    { "id": "gesamt_start",   "severity": "important",
      "value": "Die Gesamtlinie beginnt bei etwa 2,5" },
    { "id": "gesamt_end",     "severity": "important",
      "value": "endet bei etwa 2,8 Prozent des BIP" },
    { "id": "netto_negative", "severity": "important",
      "value": "Im Nettodiagramm fallen einige Jahre
       unter die Nulllinie" },
    { "id": "legend",         "severity": "important",
      "value": "Eine Legende ordnet die Farben den
       Kategorien zu" },
    { "id": "stacked",        "severity": "critical",
      "value": "gestapelte Balken" },
    { "id": "bar_color_blue", "severity": "minor",
      "value": "blau fuer Bund" },
    { "id": "bar_color_red",  "severity": "minor",
      "value": "rot fuer Laender" },
    { "id": "bar_color_yellow","severity": "minor",
      "value": "gelb fuer Gemeinden" },
    { "id": "line_color_black","severity": "minor",
      "value": "schwarze Linie fuer Gesamt" },
    { "id": "x_axis_ticks",   "severity": "minor",
      "value": "Jahresbeschriftungen auf der x-Achse" },
    { "id": "sort_order",     "severity": "minor",
      "value": "chronologisch sortiert" }
  ]
}
\end{codebox}
